% ═══════════════════════════════════════════════════════════════
% BEWERTUNGSBOGEN A1: Mündliche Prüfung Spanisch
% Rekonstruktion nach offiziellen MV-Kriterien (Handreichung KLE)
% Klasse 10 Spätbeginner - Niveau A1
% ═══════════════════════════════════════════════════════════════

\documentclass[11pt, a4paper]{article}

\usepackage[utf8]{inputenc}
\usepackage[T1]{fontenc}
\usepackage[ngerman]{babel}

% --- SCHRIFTEN ---
\usepackage{helvet}
\renewcommand{\familydefault}{\sfdefault}

\usepackage{microtype}
\usepackage[margin=1.8cm, top=2cm, bottom=1.5cm]{geometry}
\usepackage{enumitem}
\usepackage{tabularx}
\usepackage{booktabs}
\usepackage{array}
\usepackage{colortbl}
\usepackage{multirow}
\usepackage{tikz}

\usepackage{fancyhdr}
\usepackage{lastpage}

\usepackage{xcolor}
\usepackage{tcolorbox}
\tcbuselibrary{skins}

% ═══════════════════════════════════════════════════════════════
% FARBEN
% ═══════════════════════════════════════════════════════════════

\definecolor{spanischrot}{HTML}{C41E3A}
\definecolor{hellgrau}{HTML}{F5F5F5}
\definecolor{mittelgrau}{HTML}{E8E8E8}
\definecolor{dunkelgrau}{HTML}{333333}
\definecolor{rahmen}{HTML}{CCCCCC}
\definecolor{spracheblau}{HTML}{2C5AA0}
\definecolor{inhaltgruen}{HTML}{2A7A4B}

% ═══════════════════════════════════════════════════════════════
% BEFEHLE
% ═══════════════════════════════════════════════════════════════

\newcommand{\kreis}{\tikz\draw[thick] (0,0) circle (0.35em);}
\newcommand{\bewertungsfeld}{\rule{1.2cm}{0.4pt}}

\setlength{\parindent}{0pt}
\setlength{\parskip}{4pt}

% Spaltentyp für zentrierte Spalten mit fester Breite
\newcolumntype{C}[1]{>{\centering\arraybackslash}p{#1}}

% ═══════════════════════════════════════════════════════════════
% KOPF- UND FUSSZEILE
% ═══════════════════════════════════════════════════════════════

\pagestyle{fancy}
\fancyhf{}
\renewcommand{\headrulewidth}{0.5pt}
\renewcommand{\footrulewidth}{0pt}
\lhead{\footnotesize\textsc{Spanisch} · Klasse 10 · Bewertungsbogen A1}
\rhead{\footnotesize Lehrerexemplar}
\cfoot{\footnotesize\thepage\,/\,\pageref{LastPage}}

% ═══════════════════════════════════════════════════════════════
% DOKUMENT
% ═══════════════════════════════════════════════════════════════

\begin{document}

% --- TITEL ---
\begin{center}
    {\LARGE\bfseries Bewertungsbogen A1}\\[4pt]
    {\large Prueba oral -- Mündliche Prüfung Spanisch}\\[6pt]
    {\small Nach offiziellen MV-Kriterien (Handreichung KLE Sprechen)}
\end{center}

\vspace{6pt}

% --- KOPFBEREICH ---
\begin{tabularx}{\textwidth}{|X|X|}
\hline
\rowcolor{hellgrau}
\textbf{Name:} & \textbf{Datum:} \\[0.8em]
\hline
\textbf{Klasse:} \hfill 10 & \textbf{Prüfer/in:} \\[0.8em]
\hline
\end{tabularx}

\vspace{10pt}

% ═══════════════════════════════════════════════════════════════
% HINWEIS ZUR GEWICHTUNG
% ═══════════════════════════════════════════════════════════════

\begin{tcolorbox}[colback=hellgrau, colframe=hellgrau, boxrule=0pt, arc=0pt,
    left=8pt, right=8pt, top=6pt, bottom=6pt,
    borderline west={3pt}{0pt}{spanischrot}]
\textbf{Gewichtung nach MV-Handreichung:}\\[2pt]
\textcolor{spracheblau}{\textbf{Sprachliche Leistung: 60\%}} \hfill
\textcolor{inhaltgruen}{\textbf{Inhaltliche Leistung: 40\%}}\\[4pt]
{\small Bewertung: Kreuzen Sie die zutreffende Stufe an. Zwischenwerte sind möglich.}
\end{tcolorbox}

\vspace{8pt}

% ═══════════════════════════════════════════════════════════════
% TEIL A: SPRACHLICHE LEISTUNG (60%)
% ═══════════════════════════════════════════════════════════════

{\large\textcolor{spracheblau}{\textbf{A. Sprachliche Leistung}} \hfill \textbf{60\%}}

\vspace{6pt}

% --- A1: FLÜSSIGKEIT UND INTERAKTION ---
\textbf{A1. Flüssigkeit und Interaktion}

\vspace{2pt}

{\small
\begin{tabularx}{\textwidth}{|C{0.8cm}|X|C{1cm}|}
\hline
\rowcolor{mittelgrau}
\textbf{Stufe} & \textbf{Deskriptor (A1-Niveau)} & \textbf{Wertung} \\
\hline
5 & Spricht relativ flüssig in kurzen Redebeiträgen. Kann einfache Gespräche über vertraute Themen führen. Reagiert angemessen auf Gesprächspartner. & \kreis \\
\hline
4 & Spricht meist flüssig mit gelegentlichen Pausen. Kann sich an einfachen Gesprächen beteiligen. Reagiert meist angemessen. & \kreis \\
\hline
3 & Spricht mit merklichen Pausen, aber Kommunikation bleibt aufrecht. Kann auf direkte Fragen reagieren. Braucht gelegentlich Hilfe. & \kreis \\
\hline
2 & Spricht stockend mit häufigen Pausen. Reagiert verzögert. Braucht regelmäßig Unterstützung oder Wiederholung. & \kreis \\
\hline
1 & Sehr stockend, lange Pausen. Kann nur auf einfachste Fragen mit Einzelwörtern reagieren. & \kreis \\
\hline
0 & Keine zusammenhängende Kommunikation möglich. & \kreis \\
\hline
\end{tabularx}}

\vspace{8pt}

% --- A2: SPEKTRUM SPRACHLICHER MITTEL ---
\textbf{A2. Spektrum sprachlicher Mittel (Wortschatz, Strukturen)}

\vspace{2pt}

{\small
\begin{tabularx}{\textwidth}{|C{0.8cm}|X|C{1cm}|}
\hline
\rowcolor{mittelgrau}
\textbf{Stufe} & \textbf{Deskriptor (A1-Niveau)} & \textbf{Wertung} \\
\hline
5 & Verfügt über ausreichenden Wortschatz für vertraute Alltagsthemen (Familie, Wohnort, Beschreibungen). Verwendet einfache Satzstrukturen sicher (ser/estar, tener, hay, presente). & \kreis \\
\hline
4 & Grundwortschatz für die Themen vorhanden. Einfache Strukturen meist korrekt. Gelegentliche Wortsuche. & \kreis \\
\hline
3 & Begrenzter, aber funktionaler Wortschatz. Einfache Strukturen erkennbar. Häufigere Wortsuche, aber Kommunikation gelingt. & \kreis \\
\hline
2 & Sehr begrenzter Wortschatz. Nur Basisstrukturen. Häufige Verwendung von Einzelwörtern statt Sätzen. & \kreis \\
\hline
1 & Minimaler Wortschatz. Kaum erkennbare Satzstrukturen. Nur isolierte Wörter/Phrasen. & \kreis \\
\hline
0 & Kein erkennbarer Wortschatz oder Strukturen in der Zielsprache. & \kreis \\
\hline
\end{tabularx}}

\vspace{8pt}

% --- A3: KORREKTHEIT ---
\textbf{A3. Korrektheit (Grammatik, Aussprache)}

\vspace{2pt}

{\small
\begin{tabularx}{\textwidth}{|C{0.8cm}|X|C{1cm}|}
\hline
\rowcolor{mittelgrau}
\textbf{Stufe} & \textbf{Deskriptor (A1-Niveau)} & \textbf{Wertung} \\
\hline
5 & Verwendet einfache Strukturen überwiegend korrekt. Aussprache klar und verständlich. Akzente meist korrekt. Fehler beeinträchtigen Verständnis nicht. & \kreis \\
\hline
4 & Grundstrukturen meist korrekt. Aussprache gut verständlich. Einige Fehler bei komplexeren Formen, aber Verständnis gesichert. & \kreis \\
\hline
3 & Systematische Grundfehler (z.B. Verbkonjugation), aber Aussageabsicht erkennbar. Aussprache verständlich mit merklichem Akzent. & \kreis \\
\hline
2 & Häufige Fehler auch in Basisstrukturen. Aussprache teilweise schwer verständlich. Kommunikation eingeschränkt. & \kreis \\
\hline
1 & Sehr viele Fehler. Aussprache stark vom Deutschen beeinflusst. Verständigung stark beeinträchtigt. & \kreis \\
\hline
0 & Keine korrekten Strukturen erkennbar. Unverständlich. & \kreis \\
\hline
\end{tabularx}}

\vspace{10pt}

% ═══════════════════════════════════════════════════════════════
% TEIL B: INHALTLICHE LEISTUNG (40%)
% ═══════════════════════════════════════════════════════════════

{\large\textcolor{inhaltgruen}{\textbf{B. Inhaltliche Leistung}} \hfill \textbf{40\%}}

\vspace{6pt}

% --- B1: AUFGABENERFÜLLUNG ---
\textbf{B1. Aufgabenerfüllung und Themenbezug}

\vspace{2pt}

{\small
\begin{tabularx}{\textwidth}{|C{0.8cm}|X|C{1cm}|}
\hline
\rowcolor{mittelgrau}
\textbf{Stufe} & \textbf{Deskriptor (A1-Niveau)} & \textbf{Wertung} \\
\hline
5 & Erfüllt alle Teile der Aufgabe vollständig. Spricht über alle geforderten Themen (Person, Familie, Wohnort). Beantwortet Fragen inhaltlich passend. Sprachmittlung gelingt sinngemäß. & \kreis \\
\hline
4 & Erfüllt die Aufgabe größtenteils. Die meisten Themen werden angesprochen. Fragen meist passend beantwortet. Sprachmittlung im Wesentlichen gelungen. & \kreis \\
\hline
3 & Erfüllt die Aufgabe teilweise. Einige Themen fehlen oder sind nur oberflächlich. Antworten teilweise unvollständig. Sprachmittlung mit Lücken. & \kreis \\
\hline
2 & Erfüllt nur wenige Aspekte der Aufgabe. Wenige Informationen. Antworten oft nicht passend. Sprachmittlung nur ansatzweise. & \kreis \\
\hline
1 & Kaum Bezug zur Aufgabe erkennbar. Nur isolierte, nicht zusammenhängende Informationen. & \kreis \\
\hline
0 & Kein erkennbarer Bezug zur Aufgabenstellung. & \kreis \\
\hline
\end{tabularx}}

\vspace{8pt}

% --- B2: KOHÄRENZ ---
\textbf{B2. Kohärenz und Verknüpfung}

\vspace{2pt}

{\small
\begin{tabularx}{\textwidth}{|C{0.8cm}|X|C{1cm}|}
\hline
\rowcolor{mittelgrau}
\textbf{Stufe} & \textbf{Deskriptor (A1-Niveau)} & \textbf{Wertung} \\
\hline
5 & Äußerungen sind logisch aufgebaut und verknüpft. Verwendet einfache Konnektoren (y, pero, porque, también). Gedankenführung nachvollziehbar. & \kreis \\
\hline
4 & Äußerungen meist zusammenhängend. Einige Konnektoren verwendet. Gelegentliche Sprünge, aber verständlich. & \kreis \\
\hline
3 & Grundlegende Verknüpfung erkennbar. Konnektoren selten. Reihenfolge nicht immer logisch, aber Gesamtaussage verständlich. & \kreis \\
\hline
2 & Wenig Zusammenhang zwischen Äußerungen. Meist isolierte Sätze ohne Verknüpfung. & \kreis \\
\hline
1 & Nur unverbundene Einzeläußerungen. Kein erkennbarer roter Faden. & \kreis \\
\hline
0 & Keine zusammenhängenden Äußerungen. & \kreis \\
\hline
\end{tabularx}}

\newpage

% ═══════════════════════════════════════════════════════════════
% AUSWERTUNG
% ═══════════════════════════════════════════════════════════════

\begin{center}
{\Large\bfseries Auswertung}
\end{center}

\vspace{8pt}

% --- PUNKTEÜBERTRAGUNG ---
\begin{tabularx}{\textwidth}{|p{6cm}|C{2cm}|C{2cm}|X|}
\hline
\rowcolor{hellgrau}
\textbf{Kriterium} & \textbf{Stufe (0--5)} & \textbf{Gewichtung} & \textbf{Punkte} \\
\hline
\multicolumn{4}{|l|}{\cellcolor{spracheblau!15}\textbf{A. Sprachliche Leistung (60\%)}} \\
\hline
A1. Flüssigkeit und Interaktion & \bewertungsfeld & $\times$ 2 = & \bewertungsfeld \\
\hline
A2. Spektrum sprachlicher Mittel & \bewertungsfeld & $\times$ 2 = & \bewertungsfeld \\
\hline
A3. Korrektheit & \bewertungsfeld & $\times$ 2 = & \bewertungsfeld \\
\hline
\multicolumn{3}{|r|}{\textbf{Summe A (max. 30):}} & \bewertungsfeld \\
\hline
\multicolumn{4}{|l|}{\cellcolor{inhaltgruen!15}\textbf{B. Inhaltliche Leistung (40\%)}} \\
\hline
B1. Aufgabenerfüllung & \bewertungsfeld & $\times$ 2 = & \bewertungsfeld \\
\hline
B2. Kohärenz & \bewertungsfeld & $\times$ 2 = & \bewertungsfeld \\
\hline
\multicolumn{3}{|r|}{\textbf{Summe B (max. 20):}} & \bewertungsfeld \\
\hline
\end{tabularx}

\vspace{12pt}

% --- GESAMTBERECHNUNG ---
\begin{tabularx}{\textwidth}{|X|C{3cm}|}
\hline
\rowcolor{spanischrot}
\textcolor{white}{\textbf{GESAMTPUNKTZAHL}} & \textcolor{white}{\textbf{\_\_\_\_\_ / 50}} \\
\hline
\end{tabularx}

\vspace{8pt}

% --- UMRECHNUNGSTABELLE ---
\begin{center}
\textbf{Umrechnung in Notenpunkte (15-NP-Skala)}
\end{center}

\vspace{4pt}

{\small
\begin{tabularx}{\textwidth}{|*{16}{C{0.85cm}|}}
\hline
\rowcolor{hellgrau}
\textbf{NP} & 15 & 14 & 13 & 12 & 11 & 10 & 9 & 8 & 7 & 6 & 5 & 4 & 3 & 2 & 1 \\
\hline
\textbf{ab} & 49 & 48 & 47 & 45 & 42 & 40 & 37 & 33 & 30 & 27 & 23 & 20 & 17 & 13 & 10 \\
\hline
\textbf{\%} & 98 & 96 & 94 & 90 & 84 & 80 & 74 & 66 & 60 & 54 & 46 & 40 & 34 & 26 & 20 \\
\hline
\end{tabularx}}

\vspace{12pt}

% --- NOTENEINTRAG ---
\begin{tabularx}{\textwidth}{|X|C{3cm}|}
\hline
\rowcolor{spanischrot}
\textcolor{white}{\textbf{NOTENPUNKTE}} & \textcolor{white}{\textbf{\_\_\_\_\_ NP}} \\
\hline
\end{tabularx}

\vspace{16pt}

% ═══════════════════════════════════════════════════════════════
% NOTIZEN
% ═══════════════════════════════════════════════════════════════

\textbf{Notizen zur Prüfung:}

\vspace{4pt}

\begin{tabularx}{\textwidth}{|X|}
\hline
\\[4em]
\\[4em]
\\[4em]
\hline
\end{tabularx}

\vspace{12pt}

% --- UNTERSCHRIFT ---
\begin{tabularx}{\textwidth}{X X}
\rule{6cm}{0.4pt} & \rule{6cm}{0.4pt} \\
Datum & Unterschrift Prüfer/in \\
\end{tabularx}

\vspace{16pt}

% ═══════════════════════════════════════════════════════════════
% HINWEISE FÜR PRÜFER
% ═══════════════════════════════════════════════════════════════

\begin{tcolorbox}[colback=hellgrau, colframe=rahmen, boxrule=0.5pt, arc=0pt,
    left=8pt, right=8pt, top=6pt, bottom=6pt,
    title={\textbf{Hinweise zur Bewertung (A1-Niveau)}},
    coltitle=black, colbacktitle=mittelgrau]

\textbf{Tolerieren Sie bei A1-Niveau:}
\begin{itemize}[leftmargin=1.5em, itemsep=1pt, topsep=2pt]
    \item Natürliche Sprechpausen und Denkzeit
    \item Füllwörter (bueno, pues, eh...)
    \item Unvollständige Sätze
    \item Weniger komplexe Strukturen als schriftlich
    \item Wortwiederholungen
    \item Selbstkorrekturen (positiv werten!)
\end{itemize}

\vspace{4pt}

\textbf{A1-Erwartungshorizont:}
\begin{itemize}[leftmargin=1.5em, itemsep=1pt, topsep=2pt]
    \item Kann sich in einfachen, routinemäßigen Situationen verständigen
    \item Kann mit einfachen Mitteln Personen, Orte und Besitz beschreiben
    \item Kann kurze Sätze über Familie, Wohnort, Alltag bilden
    \item Kann auf direkte Fragen zu vertrauten Themen antworten
\end{itemize}

\vspace{4pt}

{\small\textit{Quelle: Rekonstruktion nach Handreichung KLE Sprechen MV (01.08.2020) und GER-Referenzrahmen A1}}
\end{tcolorbox}

\end{document}
